% =====================================================================
%  Chapitre 2 : Creation d'une application
% =====================================================================
\chapter{Créer une application Django}
\label{ch:app}

\chapitreresume{Un projet Django est composé d'applications. Dans ce chapitre,
vous créez l'application \code{catalogue}, vous découvrez le rôle de ses fichiers
et vous affichez votre première page web. Vous comprenez aussi comment une URL est
associée à une vue et comment Django renvoie une réponse au navigateur.}

\begin{objectifs}
\begin{itemize}
  \item Distinguer un projet Django d'une application Django.
  \item Créer une application avec \code{startapp}.
  \item Enregistrer l'application dans \code{INSTALLED\_APPS}.
  \item Écrire une première vue et la relier à une URL.
  \item Utiliser \code{include}, \code{app\_name} et les noms d'URL.
  \item Tester le résultat dans le navigateur.
\end{itemize}
\end{objectifs}

\section{Projet et application : quelle différence ?}

Un \textbf{projet} est l'ensemble du site : sa configuration, ses réglages, ses
URL racines. Une \textbf{application} est un module qui remplit une responsabilité
précise : catalogue, comptes, emprunts, paiements, blog, etc.

\begin{retenir}
Un projet contient plusieurs applications. Une application peut être réutilisée
dans plusieurs projets. Le projet \gesibiblio{} contiendra au minimum une
application \code{catalogue}, chargée des livres, des auteurs, des catégories et
des emprunts. Les comptes utilisateurs s'appuieront sur l'application
\code{django.contrib.auth} fournie par Django.
\end{retenir}

Au début, il est tentant de tout mettre dans une seule application géante. C'est
une mauvaise habitude. Une application doit rester cohérente : si vous devez
expliquer son rôle en une phrase, c'est bon signe. \code{catalogue} gère le
catalogue et les prêts. Plus tard, vous pourriez ajouter \code{blog},
\code{facturation} ou \code{notifications} sans toucher au catalogue.

\section{Créer l'application catalogue}

Placez-vous dans le dossier qui contient \code{manage.py}, activez
l'environnement virtuel, puis lancez :

\begin{terminal}
python manage.py startapp catalogue
\end{terminal}

Django crée le dossier \code{catalogue/} avec les fichiers suivants :

\begin{codebox}[tree]{Arborescence de l'application}
catalogue/
|-- __init__.py
|-- admin.py
|-- apps.py
|-- migrations/
|   `-- __init__.py
|-- models.py
|-- tests.py
`-- views.py
\end{codebox}

\begin{note}
\code{startapp} ne crée pas de fichier \code{urls.py}. C'est volontaire : toutes
les applications n'ont pas besoin d'URL propres. Nous allons le créer à la main
pour garder un routage clair.
\end{note}

\section{Le rôle de chaque fichier de l'application}

\begin{table}[htbp]
\centering
\small
\begin{tabularx}{\textwidth}{L{3.6cm} X}
\toprule
\textbf{Fichier} & \textbf{Rôle} \\
\midrule
\code{\_\_init\_\_.py} & Transforme le dossier en paquet Python. Il reste vide. \\
\code{admin.py} & Enregistre les modèles dans l'interface d'administration. \\
\code{apps.py} & Configuration de l'application : nom, libellé, éventuels
branchements au démarrage. \\
\code{models.py} & Définit les modèles de données, c'est-à-dire les tables et
leurs relations. \\
\code{views.py} & Contient les vues, les fonctions qui reçoivent les requêtes et
renvoient les réponses. \\
\code{tests.py} & Contient les tests automatisés de l'application. \\
\code{migrations/} & Historique des modifications du schéma de base de données. \\
\code{urls.py} & À créer : table de routage propre à l'application. \\
\bottomrule
\end{tabularx}
\caption{Rôle des fichiers d'une application Django.}
\end{table}

\subsection{apps.py}

\begin{codebox}[python]{catalogue/apps.py}
from django.apps import AppConfig


class CatalogueConfig(AppConfig):
    default_auto_field = "django.db.models.BigAutoField"
    name = "catalogue"
    verbose_name = "Catalogue"
\end{codebox}

\code{AppConfig} décrit l'application pour Django. Le champ \code{name} doit
correspondre exactement au nom du paquet Python. \code{verbose\_name} est le nom
affiché dans l'administration.

\section{Enregistrer l'application}

Une application n'est active que si elle figure dans la liste des applications
installées. Ouvrez le fichier de réglages du projet et ajoutez la ligne suivante
à la fin de cette liste :

\begin{codebox}[python]{gesibiblio/settings.py, extrait}
INSTALLED_APPS = [
    "django.contrib.admin",
    "django.contrib.auth",
    "django.contrib.contenttypes",
    "django.contrib.sessions",
    "django.contrib.messages",
    "django.contrib.staticfiles",
    "catalogue.apps.CatalogueConfig",
]
\end{codebox}

Vous pouvez aussi écrire simplement \code{"catalogue"}. Django retrouvera la
classe \code{CatalogueConfig} grâce à la convention de nommage. La forme complète
est plus explicite et permet de personnaliser le libellé.

\begin{attention}
Si vous oubliez d'enregistrer l'application, Django ne verra ni ses modèles, ni
ses gabarits, ni ses fichiers statiques. Les erreurs du type \code{No installed
app with label 'catalogue'} viennent presque toujours de là.
\end{attention}

\section{Première vue}

Une vue est une fonction Python qui reçoit un objet \code{request} et renvoie un
objet \code{response}. Pour commencer simplement, nous allons renvoyer du HTML
brut avec \code{HttpResponse}.

\begin{codebox}[python]{catalogue/views.py}
from django.http import HttpResponse


def accueil(request):
    return HttpResponse(
        "<h1>Bienvenue dans GesiBiblio</h1>"
        "<p>La bibliothèque ouvre bientôt ses portes.</p>"
    )
\end{codebox}

\begin{description}[leftmargin=3.1cm,style=nextline]
  \item[\code{request}] Représente la requête HTTP. Il contient la méthode, l'URL,
  les paramètres, l'utilisateur connecté, les cookies, etc. Même si nous ne
  l'utilisons pas encore, il doit toujours être présent dans la signature.
  \item[\code{HttpResponse}] Représente la réponse HTTP. Le premier argument est
  le corps de la réponse, ici une chaîne HTML. Django ajoute automatiquement les
  en-têtes et le code de statut \code{200}.
\end{description}

\begin{note}
Écrire du HTML dans une vue est utile pour comprendre le mécanisme, mais ce n'est
pas une bonne pratique. Dès le chapitre suivant, nous utiliserons des gabarits
pour séparer la présentation de la logique.
\end{note}

\section{Relier une URL à la vue}

Créez le fichier \code{catalogue/urls.py} :

\begin{codebox}[python]{catalogue/urls.py}
from django.urls import path
from . import views

app_name = "catalogue"

urlpatterns = [
    path("", views.accueil, name="accueil"),
]
\end{codebox}

Puis modifiez le fichier \code{gesibiblio/urls.py} pour inclure ces URL :

\begin{codebox}[python]{gesibiblio/urls.py}
from django.contrib import admin
from django.urls import include, path

urlpatterns = [
    path("admin/", admin.site.urls),
    path("", include("catalogue.urls")),
]
\end{codebox}

\subsection{Comprendre path et include}

\begin{description}[leftmargin=3.1cm,style=nextline]
  \item[\code{path("", vue, name="...")}] Associe une URL à une vue. Ici, la
  chaîne vide correspond à la racine du site. Le nom \code{accueil} permet de
  construire l'URL dans les gabarits et dans le code Python.
  \item[\code{include("catalogue.urls")}] Délègue le routage à l'application. Toute
  URL non consommée par les routes précédentes est cherchée dans
  \code{catalogue/urls.py}. C'est ce qui rend les applications autonomes.
  \item[\code{app\_name = "catalogue"}] Crée un espace de noms. Le nom complet de
  la route devient \code{catalogue:accueil}. Deux applications peuvent ainsi avoir
  une route \code{accueil} sans se marcher dessus.
\end{description}

\begin{bonne}
Nommez toujours vos routes. Une URL codée en dur dans un gabarit casse dès que le
chemin change. \lstinline|{% url 'catalogue:accueil' %}| résiste aux modifications
et signale les erreurs à la compilation des gabarits.
\end{bonne}

\section{Tester la page}

Relancez le serveur :

\begin{terminal}
python manage.py runserver
\end{terminal}

Ouvrez \code{http://127.0.0.1:8000/}. Vous devez voir :

\begin{quote}
\og Bienvenue dans GesiBiblio \fg\\
\og La bibliothèque ouvre bientôt ses portes. \fg
\end{quote}

Le navigateur affiche le HTML brut, sans mise en forme. C'est normal : nous
n'avons pas encore de feuille de style.

\begin{attention}
Si vous obtenez une erreur \code{TemplateDoesNotExist}, c'est que vous avez déjà
ajouté un appel à \code{render} sans créer le gabarit. Si vous obtenez
\code{Page not found}, vérifiez l'ordre des \code{path} et le contenu de
\code{urlpatterns}. Si vous obtenez \code{ModuleNotFoundError}, vérifiez que
l'application est bien dans \code{INSTALLED\_APPS} et que le dossier contient
\code{\_\_init\_\_.py}.
\end{attention}

\section{Ajouter une deuxième page}

Pour voir la mécanique en action, ajoutons une page \og À propos \fg.

\begin{codebox}[python]{catalogue/views.py}
from django.http import HttpResponse


def accueil(request):
    return HttpResponse(
        "<h1>Bienvenue dans GesiBiblio</h1>"
        "<p>La bibliothèque ouvre bientôt ses portes.</p>"
    )


def a_propos(request):
    return HttpResponse(
        "<h1>À propos de GesiBiblio</h1>"
        "<p>Un projet de bibliothèque réalisé avec Django.</p>"
    )
\end{codebox}

\begin{codebox}[python]{catalogue/urls.py}
from django.urls import path
from . import views

app_name = "catalogue"

urlpatterns = [
    path("", views.accueil, name="accueil"),
    path("a-propos/", views.a_propos, name="a_propos"),
]
\end{codebox}

Visitez \code{http://127.0.0.1:8000/a-propos/}. La page s'affiche.

\begin{note}
L'URL se termine par une barre oblique. Django peut ajouter automatiquement cette
barre grâce au réglage \code{APPEND\_SLASH}, activé par défaut. Prenez l'habitude
d'écrire des URL cohérentes, avec une barre finale.
\end{note}

\section{Le cycle requête réponse}

Voici ce qui se passe lorsque vous demandez \code{/a-propos/} :

\begin{table}[htbp]
\centering
\small
\begin{tabularx}{\textwidth}{C{1.1cm} L{4.1cm} X}
\toprule
\textbf{\#} & \textbf{Composant} & \textbf{Action} \\
\midrule
1 & Navigateur & Envoie une requête \code{GET /a-propos/}. \\
2 & Middlewares & Ajoutent la session, l'utilisateur, la protection CSRF. \\
3 & \code{gesibiblio/urls.py} & Ne reconnaît pas \code{admin/}, délègue à
\code{catalogue.urls}. \\
4 & \code{catalogue/urls.py} & Reconnaît \code{a-propos/} et appelle
\code{a\_propos}. \\
5 & Vue & Reçoit \code{request}, prépare la réponse. \\
6 & \code{HttpResponse} & Renvoie le HTML avec le statut \code{200}. \\
7 & Navigateur & Affiche le corps de la réponse. \\
\bottomrule
\end{tabularx}
\caption{Trajet d'une requête vers la page \og À propos \fg.}
\end{table}

\begin{retenir}
Une URL n'est pas un fichier. C'est une route décrite dans \code{urls.py}. Le
serveur ne cherche pas \code{a-propos.html} sur le disque : il appelle la fonction
\code{a\_propos}, qui construit la réponse.
\end{retenir}

\section{Bonnes pratiques}

\begin{bonne}
\begin{itemize}
  \item Une application, une responsabilité claire.
  \item Un fichier \code{urls.py} par application, inclus depuis le routeur racine.
  \item Un \code{app\_name} et des noms de routes explicites.
  \item Des URL en minuscules, avec des tirets, et une barre finale.
  \item Des vues courtes : la logique métier lourde va dans les modèles, les
        formulaires ou des services dédiés.
  \item Pas de HTML dans les vues : utilisez des gabarits dès que possible.
\end{itemize}
\end{bonne}

\begin{piege}
Deux erreurs classiques débutent par la même confusion : \code{include} ne
fonctionne que si le module cible expose une variable \code{urlpatterns}, et
\code{app\_name} doit être identique au préfixe utilisé dans
\lstinline|{% url 'catalogue:...' %}|. Vérifiez ces deux points avant de chercher
ailleurs.
\end{piege}

\section{Exercices}

\begin{exercice}[ : page de contact]
Ajoutez une page \code{/contact/} qui affiche une adresse électronique et les
horaires d'ouverture de la bibliothèque. Donnez-lui le nom de route
\code{contact}.
\end{exercice}

\begin{correction}
\begin{codebox}[python]{catalogue/views.py, extrait}
def contact(request):
    return HttpResponse(
        "<h1>Contact</h1>"
        "<p>Écrivez-nous à contact@gesibiblio.example.</p>"
        "<p>Ouvert du lundi au samedi, de 9 h à 18 h.</p>"
    )
\end{codebox}

\begin{codebox}[python]{catalogue/urls.py, extrait}
path("contact/", views.contact, name="contact"),
\end{codebox}

Vérifiez ensuite \code{http://127.0.0.1:8000/contact/}.
\end{correction}

\begin{exercice}[ : lien entre les pages]
Ajoutez un lien vers la page \og À propos \fg{} sur la page d'accueil, puis un lien
de retour vers l'accueil sur la page \og À propos \fg. Utilisez les noms de routes
dans le HTML.
\end{exercice}

\begin{correction}
Dans une vue qui renvoie du HTML brut, on peut utiliser \code{reverse} :

\begin{codebox}[python]{catalogue/views.py, extrait}
from django.urls import reverse


def accueil(request):
    lien = reverse("catalogue:a_propos")
    return HttpResponse(
        "<h1>Bienvenue dans GesiBiblio</h1>"
        f'<p><a href="{lien}">En savoir plus</a></p>'
    )
\end{codebox}

Dans un vrai gabarit, la balise \lstinline|{% url 'catalogue:a_propos' %}| est
plus lisible. Nous la verrons au chapitre suivant.
\end{correction}

\begin{exercice}[ : comprendre les noms]
Expliquez la différence entre \code{path("livres/", ...)} et
\code{path("livres/", ..., name="livre\_liste")}. Pourquoi le nom est-il utile ?
\end{exercice}

\begin{correction}
Le premier argument est le motif d'URL reconnu par le navigateur. Le nom
\code{livre\_liste} est un identifiant interne. Il permet de construire l'URL avec
\code{reverse("catalogue:livre\_liste")} ou
\lstinline|{% url 'catalogue:livre_liste' %}|. Si le chemin change demain, seul
\code{urls.py} est modifié : les gabarits et le code continuent de fonctionner.
\end{correction}

\begin{retenir}
\begin{itemize}
  \item Un projet contient des applications ; une application contient des vues,
        des modèles, des gabarits et des URL.
  \item \code{startapp} crée la structure de l'application ; il faut l'enregistrer
        dans \code{INSTALLED\_APPS}.
  \item Une vue reçoit \code{request} et renvoie une réponse.
  \item \code{include} délègue le routage à une application.
  \item \code{app\_name} et \code{name} permettent de nommer les routes.
  \item Le HTML brut est utile pour apprendre, mais les gabarits arrivent vite.
\end{itemize}
\end{retenir}
